% GENERATED FILE. Edit canonical lessons, then run npm run generate:books.
% canonical-source-hash: fnv1a64:52c0d2e5ab6ae295
% canonical-lessons: ES-C622-normativa, ES-C622-reglamento, ES-C622-prohibicion, ES-C622-obligar, ES-C622-obligacion

\chapter{Las reglas}
\label{ch:es-las-reglas}
\hlchaptermodality{622}

\begin{chapteropening}
\textbf{By the end of this chapter:} \emph{I can read the rules of a place and say what is obligatory.}

Complete the last lesson of chapter 622: I can read the rules of a place and say what is obligatory.
\end{chapteropening}

\section[la normativa]{la normativa — the set of rules, the regulations}

\label{lesson:ES-C622-normativa}



\begin{quote}
Before the new one: say the Spanish for exempt, then the Spanish for forbidden.
\end{quote}

\subsection*{You'll want to know: la normativa}
\textbf{la normativa} — \textquotedblleft{}the set of rules, the regulations\textquotedblright{}.

\begin{quote}
La nueva normativa empieza en enero.
\end{quote}

\begin{quote}
\emph{The new rules start in January.}
\end{quote}

Rules made by an authority for a whole area.

\begin{cousinweb}
\emph{normativa} is built on \emph{norma}, which comes from Latin \emph{norma} 'a carpenter's square, a rule' (EtymDB).
\end{cousinweb}

\begin{grammarlens}[title={What you've built}]
One more word for telling, explaining and describing at B1.
\end{grammarlens}

\subsection*{Your turn}
\begin{itemize}
\raggedright
  \item \emph{Say it:} \emph{la normativa}
  \item \emph{Say it:} \emph{la normativa}, once more
  \item \emph{Say it:} a sentence of your own with \emph{la normativa}
  \item \emph{Recall:} say the Spanish for exempt, then the Spanish for forbidden, then say \emph{la normativa} again
\end{itemize}

\begin{tcolorbox}[breakable,colback=teal!4,colframe=teal!35!black,title={Before you move on}]
What does la normativa mean? (\textquotedblleft{}The set of rules, the regulations\textquotedblright{}.) Say it once more.
\end{tcolorbox}
\section[el reglamento]{el reglamento — the regulations, the rules}

\label{lesson:ES-C622-reglamento}



\begin{quote}
Before the new one: say the Spanish for forbidden, then the Spanish for the set of rules.
\end{quote}

\subsection*{You'll want to know: el reglamento}
\textbf{el reglamento} — \textquotedblleft{}the regulations, the rules\textquotedblright{}.

\begin{quote}
Lee el reglamento antes de empezar.
\end{quote}

\begin{quote}
\emph{Read the regulations before you start.}
\end{quote}

The written set of rules of a club, a building or a game.

\begin{cousinweb}
\emph{reglamento} is built on \emph{reglar}, a learned borrowing from Latin \emph{regulare} 'to direct', from \emph{regula} 'rule' (EtymDB).
\end{cousinweb}

\begin{grammarlens}[title={What you've built}]
One more word for telling, explaining and describing at B1.
\end{grammarlens}

\subsection*{Your turn}
\begin{itemize}
\raggedright
  \item \emph{Say it:} \emph{el reglamento}
  \item \emph{Say it:} \emph{el reglamento}, once more
  \item \emph{Say it:} a sentence of your own with \emph{el reglamento}
  \item \emph{Recall:} say the Spanish for forbidden, then the Spanish for the set of rules, then say \emph{el reglamento} again
\end{itemize}

\begin{tcolorbox}[breakable,colback=teal!4,colframe=teal!35!black,title={Before you move on}]
What does el reglamento mean? (\textquotedblleft{}The regulations, the rules\textquotedblright{}.) Say it once more.
\end{tcolorbox}
\section[la prohibición]{la prohibición — the ban, the prohibition}

\label{lesson:ES-C622-prohibicion}



\begin{quote}
Before the new one: say the Spanish for the set of rules, then the Spanish for the regulations.
\end{quote}

\subsection*{You'll want to know: la prohibición}
\textbf{la prohibición} — \textquotedblleft{}the ban, the prohibition\textquotedblright{}.

\begin{quote}
Hay una prohibición de fumar.
\end{quote}

\begin{quote}
\emph{There is a smoking ban.}
\end{quote}

Built on \emph{prohibir}.

\begin{cousinweb}
\emph{prohibición} is a learned borrowing from Latin \emph{prohibitio} 'a forbidding', built on \emph{prohibere} 'to hold back', from \emph{habere} 'to hold' (Lewis \& Short).
\end{cousinweb}

\begin{grammarlens}[title={What you've built}]
One more word for telling, explaining and describing at B1.
\end{grammarlens}

\subsection*{Your turn}
\begin{itemize}
\raggedright
  \item \emph{Say it:} \emph{la prohibición}
  \item \emph{Say it:} \emph{la prohibición}, once more
  \item \emph{Say it:} a sentence of your own with \emph{la prohibición}
  \item \emph{Recall:} say the Spanish for the set of rules, then the Spanish for the regulations, then say \emph{la prohibición} again
\end{itemize}

\begin{tcolorbox}[breakable,colback=teal!4,colframe=teal!35!black,title={Before you move on}]
What does la prohibición mean? (\textquotedblleft{}The ban, the prohibition\textquotedblright{}.) Say it once more.
\end{tcolorbox}
\section[obligar]{obligar — to force, to oblige}

\label{lesson:ES-C622-obligar}



\begin{quote}
Before the new one: say the Spanish for the regulations, then the Spanish for the ban.
\end{quote}

\subsection*{You'll want to know: obligar}
\textbf{obligar} — \textquotedblleft{}to force, to oblige\textquotedblright{}.

\begin{quote}
Nadie te obliga a venir.
\end{quote}

\begin{quote}
\emph{Nobody is forcing you to come.}
\end{quote}

\emph{Obligar a alguien a} + infinitive. Spelling: \emph{obligué}.

\begin{cousinweb}
\emph{obligar} is a learned borrowing from Latin \emph{obligare} 'to bind, to put under a duty', built on \emph{ligare} 'to tie' (EtymDB).
\end{cousinweb}

\begin{grammarlens}[title={What you've built}]
One more word for telling, explaining and describing at B1.
\end{grammarlens}

\subsection*{Your turn}
\begin{itemize}
\raggedright
  \item \emph{Say it:} \emph{obligar}
  \item \emph{Say it:} \emph{obligar}, once more
  \item \emph{Say it:} a sentence of your own with \emph{obligar}
  \item \emph{Recall:} say the Spanish for the regulations, then the Spanish for the ban, then say \emph{obligar} again
\end{itemize}

\begin{tcolorbox}[breakable,colback=teal!4,colframe=teal!35!black,title={Before you move on}]
What does obligar mean? (\textquotedblleft{}To force, to oblige\textquotedblright{}.) Say it once more.
\end{tcolorbox}
\section[la obligación]{la obligación — the obligation, the duty}

\label{lesson:ES-C622-obligacion}



\begin{quote}
Before the new one: say the Spanish for the ban, then the Spanish for to force.
\end{quote}

\subsection*{You'll want to know: la obligación}
\textbf{la obligación} — \textquotedblleft{}the obligation, the duty\textquotedblright{}.

\begin{quote}
Votar es un derecho, no una obligación.
\end{quote}

\begin{quote}
\emph{Voting is a right, not an obligation.}
\end{quote}

\emph{Tener la obligación de} + infinitive.

\begin{cousinweb}
\emph{obligación} is a learned borrowing from Latin \emph{obligatio} 'a binding, a legal tie', built on \emph{obligare} 'to bind' (EtymDB).
\end{cousinweb}

\begin{grammarlens}[title={What you've built}]
One more word for telling, explaining and describing at B1.
\end{grammarlens}

\subsection*{Your turn}
\begin{itemize}
\raggedright
  \item \emph{Say it:} \emph{la obligación}
  \item \emph{Say it:} \emph{la obligación}, once more
  \item \emph{Say it:} a sentence of your own with \emph{la obligación}
  \item \emph{Recall:} say the Spanish for the ban, then the Spanish for to force, then say \emph{la obligación} again
\end{itemize}

\begin{tcolorbox}[breakable,colback=teal!4,colframe=teal!35!black,title={Before you move on}]
What does la obligación mean? (\textquotedblleft{}The obligation, the duty\textquotedblright{}.) Say it once more.
\end{tcolorbox}
